\section{The Klein multiplicity structure}
\label{sec:klein}

Let \(X\) be the Klein cubic threefold and \(J(X)\) its intermediate
Jacobian.  With the standard Tate-twist convention,
\[
 H^1(J(X),\Q)\simeq H^3(X,\Q)(1).
\]
Roulleau's period calculation \cite{RoulleauKlein} gives an unpolarized complex-abelian-variety
isomorphism
\[
 J(X)\cong E^5,\qquad
 \End_{\mathbf C}^0(E)=\Q(\sqrt{-11}).
\]

The ten-dimensional rational representation has no two-dimensional
\(\PSL_2(11)\)- or \(A_5\)-stable subspace.  Thus the orientation carrier
cannot be an equivariant elliptic subvariety.  Upon restriction to either
icosahedral parent, however,
\[
 H^1(J(X),\Q)|_{A_5}\simeq W_5\otimes U,
\qquad
 U=\Hom_{A_5}(W_5,H^1(J(X),\Q)),
\]
where \(W_5\) is the rational irreducible five-space and \(U\) is a
two-dimensional weight-one Hodge structure.

For the two parents \(A_5^\pm\), set
\[
 C_\pm=\End_{A_5^\pm}(H^1(J(X),\Q)).
\]
The expected exact relations are
\[
 C_\pm\simeq M_2(\Q),\qquad
 C_+\cap C_-\simeq\Q(\sqrt{-11}).
\]
The first field records the relative icosahedral position only if it is
extracted by a canonical mixed invariant.  Arbitrary rank-one idempotents
in \(M_2(\Q)\) are not canonical.

\paragraph{Strong-form target.}
Using the principal polarization and its Rosati trace form, construct an
invariant of the ordered pair \((C_+,C_-)\) whose discriminant class is
\(5\).  This would recover \(\Q(\sqrt5)\) transversely to the CM field
\(\Q(\sqrt{-11})\).  The exact rational representation, polarization,
commutants, and invariant are assigned to \claimid{KL-1},
\claimid{KL-2}, and \claimid{KL-3}; they remain open under C654.

The Fano surface also contains \(55\) canonical genus-two curves
\cite{RoulleauCurves} indexed by
the involutions of \(\PSL_2(11)\).  Their intersection lattice has rank
\(25\) and index \(2\) in the Néron--Severi group.  Whether that saturation
detects the finite orientation bit is a second-stage question, not evidence
for the strong-form target; it is recorded as \claimid{KL-4}.
